\begin{lemma}[Every copied row reaches atoms]
If $h$ is bad for $\mathsf{PowerRel}(r)$, then for every $x$ and $n$ there
is a finite depth $k$ such that
$\mathsf{PowerCopiedTerminalAt}_r(h,x,n,k)$ holds.
\end{lemma}
\begin{proof}
Write
\[
 L_{m,j}=\mathsf{PowerCopiedLeft}_r(h,x,m,j)
 \quad\text{and}\quad
 P_j=(L_{n,j},L_{n+1,j}).
\]
For every $j$, \noderef{PowerCopiedFailureMove} gives both
\[
 \neg\mathsf{PowerRel}(r)(L_{n,j},L_{n+1,j})
 \quad\text{and}\quad
 \mathsf{PowerMove}(L_{n,j},L_{n,j+1}),
\]
and, applying the same statement with row $n+1$, it gives
\[
 \mathsf{PowerMove}(L_{n+1,j},L_{n+1,j+1}).
\]
Thus every $P_j$ is still a failed adjacent copied pair, while the two
displayed move relations describe its next pair.

We claim that whenever $P_j$ is not a pair of atoms,
\[
 \mathsf{PowerGameDescent}(P_{j+1},P_j)                       \tag{*}
\]
holds.  If $L_{n,j}$ is a node, the definition of
\noderef{PowerMove} says that the first move replaces it by one of its
children.  Hence
\(\mathsf{PowerChild}(L_{n,j+1},L_{n,j})\), and the first disjunct in
\noderef{PowerGameDescent} proves $(*)$.  If $L_{n,j}$ is already an atom,
the same definition of \noderef{PowerMove} says
\(L_{n,j+1}=L_{n,j}\).  Since $P_j$ is not a pair of atoms, its second
coordinate $L_{n+1,j}$ is then a node.  The second move relation therefore
gives
\(\mathsf{PowerChild}(L_{n+1,j+1},L_{n+1,j})\), so the second disjunct in
\noderef{PowerGameDescent} proves $(*)$.

By \noderef{PowerGameDescentWellFounded}, $P_0$ is accessible for
\(\mathsf{PowerGameDescent}\).  We use accessibility induction along the
fixed sequence $(P_j)_j$.  At an accessible $P_j$, if it is a pair of
atoms, we stop at depth $j$.  Otherwise $(*)$ makes $P_{j+1}$ an accessible
predecessor, so the induction hypothesis supplies a later index at which
the pair is made of two atoms.  Therefore an infinite nonterminal
continuation is impossible, and some finite $k$ has both $L_{n,k}$ and
$L_{n+1,k}$ atoms.  By the definition in
\noderef{PowerCopiedTerminalAt}, choosing those two atom values as $q$ and
$q'$ gives
\(\mathsf{PowerCopiedTerminalAt}_r(h,x,n,k)\), as required.
\end{proof}
